% tcy-fontchange.tex -- 縦中横の内部フォント変更（\gtfamily 等）と
% 境界（前後のJFMグルー二重挿入なし・kinsoku 保持）のゲート。
% 契約:
%   * \tatechuyoko{...\gtfamily 12} のような内部の書体変更が
%     1 回だけの組版で反映される（工場は body を一度だけ組む）。
%   * 前後の和文との境界グルーは二重に挿入されない（before-glue は
%     \jlreq@calc@jfmglue@cp が 1 回、box→後続字の境界は rules の
%     評価器が 1 回）。
%   * 非star形の yoko での使用は ClassError（別 fixture tcy-star-error）。
\documentclass[pdflatex,tate,paper=a5]{jlreq}
\pagestyle{empty}
\makeatletter
\begin{document}
% --- TCY-GT: 内部 \gtfamily 変更（body は一度だけ組まれbox寸法に反映）
\newbox\tcybox
\setbox\tcybox=\hbox{\tatechuyoko{\gtfamily 12}}%
\ifdim\wd\tcybox=\dimexpr1\jlreq@jareg\relax
  \typeout{TATE-PASS: TCY-GT (wd=\the\wd\tcybox)}%
\else
  \typeout{TATE-FAIL: TCY-GT (wd=\the\wd\tcybox)}%
\fi
% --- TCY-BF: 内部 \bfseries
\setbox\tcybox=\hbox{\tatechuyoko{\bfseries 34}}%
\ifdim\wd\tcybox=\dimexpr1\jlreq@jareg\relax
  \typeout{TATE-PASS: TCY-BF (wd=\the\wd\tcybox)}%
\else
  \typeout{TATE-FAIL: TCY-BF (wd=\the\wd\tcybox)}%
\fi
% --- TCY-BOUNDARY: 直前和文→縦中横→直後和文の境界で追加グルーは
% uplatex と同じくほぼ 0（阿→box と box→阿）。実測: セル数を揃えて
% 「阿伊\tatechuyoko{12}阿伊阿」（2+tcy+2 = 5 セル相当）と
% 「阿伊阿阿伊阿」（素の漢字 5 セル）の行長差。許容 0.5pt
% （境界グルーが 0 でない場合に破綻するのは 1zw=10pt 級のズレ）。
\newdimen\lenwithtcy \newdimen\lenplain
\setbox\tcybox=\hbox{阿伊\tatechuyoko{12}阿伊阿}%
\lenwithtcy=\wd\tcybox
\setbox\tcybox=\hbox{阿伊阿阿伊阿}%
\lenplain=\wd\tcybox
\ifdim\dimexpr\lenwithtcy-\lenplain\relax<0.5pt\relax
  \ifdim\dimexpr\lenplain-\lenwithtcy\relax<0.5pt\relax
    \typeout{TATE-PASS: TCY-BOUNDARY (delta=\the\dimexpr\lenwithtcy-\lenplain\relax)}%
  \else
    \typeout{TATE-FAIL: TCY-BOUNDARY (delta=\the\dimexpr\lenwithtcy-\lenplain\relax)}%
  \fi
\else
  \typeout{TATE-FAIL: TCY-BOUNDARY (delta=\the\dimexpr\lenwithtcy-\lenplain\relax, expected within 0.5pt of 0)}%
\fi
% --- 実レンダリング
阿伊\tatechuyoko{\gtfamily 12}阿伊\tatechuyoko{\bfseries 56}阿伊
縦中横と、句読点。\tatechuyoko{12}「引用」の禁則も保持される。
\end{document}
